% Options for packages loaded elsewhere
\PassOptionsToPackage{unicode}{hyperref}
\PassOptionsToPackage{hyphens}{url}
\documentclass[
]{article}
\usepackage{xcolor}
\usepackage{amsmath,amssymb}
\setcounter{secnumdepth}{-\maxdimen} % remove section numbering
\usepackage{iftex}
\ifPDFTeX
  \usepackage[T1]{fontenc}
  \usepackage[utf8]{inputenc}
  \usepackage{textcomp} % provide euro and other symbols
\else % if luatex or xetex
  \usepackage{unicode-math} % this also loads fontspec
  \defaultfontfeatures{Scale=MatchLowercase}
  \defaultfontfeatures[\rmfamily]{Ligatures=TeX,Scale=1}
\fi
\usepackage{lmodern}
\ifPDFTeX\else
  % xetex/luatex font selection
\fi
% Use upquote if available, for straight quotes in verbatim environments
\IfFileExists{upquote.sty}{\usepackage{upquote}}{}
\IfFileExists{microtype.sty}{% use microtype if available
  \usepackage[]{microtype}
  \UseMicrotypeSet[protrusion]{basicmath} % disable protrusion for tt fonts
}{}
\makeatletter
\@ifundefined{KOMAClassName}{% if non-KOMA class
  \IfFileExists{parskip.sty}{%
    \usepackage{parskip}
  }{% else
    \setlength{\parindent}{0pt}
    \setlength{\parskip}{6pt plus 2pt minus 1pt}}
}{% if KOMA class
  \KOMAoptions{parskip=half}}
\makeatother
\setlength{\emergencystretch}{3em} % prevent overfull lines
\providecommand{\tightlist}{%
  \setlength{\itemsep}{0pt}\setlength{\parskip}{0pt}}
\usepackage{bookmark}
\IfFileExists{xurl.sty}{\usepackage{xurl}}{} % add URL line breaks if available
\urlstyle{same}
\hypersetup{
  hidelinks,
  pdfcreator={LaTeX via pandoc}}

\author{}
\date{}

\begin{document}

\section{Proposal: Cloud-Native Multimodal AI Study Assistant
(GCP)}\label{proposal-cloud-native-multimodal-ai-study-assistant-gcp}

\subsection{1. Problem Statement}\label{problem-statement}

Students and early-career professionals often struggle to combine
planning, research, and automation in one place when learning technical
topics. They switch between search engines, notes apps, shell tools, and
calendars, which adds friction and reduces focus. This project addresses
that problem by building a cloud-hosted AI assistant that can answer
questions, summarize findings, automate simple tasks, and support voice
interaction.

\subsection{2. Intended Users}\label{intended-users}

\begin{itemize}
\tightlist
\item
  Primary users: university students learning cloud and AI topics.
\item
  Secondary users: professionals who need quick decision support for
  scheduling, research, or planning.
\item
  Personal use: self-study assistant for structured weekly goals and
  actionable task breakdowns.
\end{itemize}

\subsection{3. Cloud Platform and Planned
Services}\label{cloud-platform-and-planned-services}

Platform: Google Cloud Platform (GCP)

Planned cloud-native components: - Cloud Run: hosts the assistant
backend API. - Cloud Functions (2nd gen): external tool endpoint for
automation actions. - Firestore: lightweight persistent storage for
assistant notes/facts. - Cloud Speech-to-Text: transcribes user voice
input. - Cloud Text-to-Speech: converts assistant responses to audio. -
Artifact Registry + Cloud Build: image storage and CI/CD deployment
pipeline.

Cross-cloud/hybrid extension: - Hybrid mode with local client + cloud
backend (local user interface calling GCP-hosted API). - Optional
fallback LLM provider integration for interoperability testing.

\subsection{4. Type of AI Assistant}\label{type-of-ai-assistant}

This project will implement a multimodal cloud AI assistant with: - Text
chat for Q and A, planning, and summarization. - Voice input/output for
accessibility and hands-free use. - Tool-calling support for web search,
database lookup, cloud function automation, and sandboxed shell
commands. - Practical use case: generating weekly study plans,
summarizing recent topic updates, and running automation actions (for
example, getting UTC time).

\subsection{Proposed Deliverables}\label{proposed-deliverables}

\begin{itemize}
\tightlist
\item
  Deployed assistant endpoint running on Cloud Run.
\item
  At least one AI model/API: Gemini LLM, plus speech services.
\item
  Cloud-native architecture with serverless services and storage.
\item
  Evaluation results using defined test cases.
\item
  Final report with architecture, data flow, metrics, limitations, and
  future work.
\end{itemize}

\subsection{Security and Privacy
Approach}\label{security-and-privacy-approach}

Identified risks: - API key leakage from source code or logs. - Command
injection through tool-calling shell interface.

Mitigations: - API keys stored in environment variables (or Secret
Manager in production), never hard-coded. - Input sanitization and
strict allow-list policy for shell commands.

\subsection{Success Criteria}\label{success-criteria}

\begin{itemize}
\tightlist
\item
  Assistant can accept input and return useful AI-generated output.
\item
  At least one practical scenario is completed end-to-end.
\item
  Evaluation average score \textgreater= 0.70 on rubric-based test set.
\item
  Demonstrated multimodal value over text-only baseline.
\end{itemize}

\end{document}
